%% Dokumentklasse DHBW Bachelor
\documentclass[
    12pt,              % Schriftgröße
    a4paper,           % Papierformat
    oneside,           % Einseitiger Druck
    openany,           % Kapitel können auf jeder Seite beginnen
    parskip=half,      % Absatzabstand
    headings=normal    % Überschriftenformat
]{scrbook}

%% ============================================================================
%% METADATEN - BITTE ANPASSEN
%% ============================================================================

% Dokumenttitel
\newcommand{\doctitle}{KI-unterstützte Erstellung von ETL-Prozessen}
\newcommand{\docsubtitle}{Eine experimentelle Evaluierung}

% Autor
\newcommand{\authorname}{Luis Geiger}
\newcommand{\matnr}{9144849}
\newcommand{\course}{Informatik}

% Betreuer
\newcommand{\supervisor}{Jens Vierling}
\newcommand{\company}{Mercedes-Benz Tech Innovation GmbH}
\newcommand{\companyaddress}{Gropiusplatz 10, 70563 Stuttgart}

% Dauer
\newcommand{\startdate}{14.06.2026}
\newcommand{\dateend}{07.09.2026}
\newcommand{\submitdate}{\today}

%% ============================================================================
%% PAKETE UND KONFIGURATION
%% ============================================================================

% Sprache und Encoding
\usepackage[ngerman]{babel}
\usepackage[utf8]{inputenc}
\usepackage[T1]{fontenc}
\usepackage{lmodern}

% Seiten und Ränder
\usepackage[
    left=2.5cm,
    right=2.5cm,
    top=2.5cm,
    bottom=2.5cm,
    headheight=1.5cm,
    footskip=1.5cm
]{geometry}

% Zitationen und Bibliografie
\usepackage[backend=biber, style=authoryear, sorting=nyt]{biblatex}
\addbibresource{resources/bibliography.bib}

% Grafiken und Tabellen
\usepackage{graphicx}
\usepackage{xcolor}
\usepackage{colortbl}
\usepackage{tabularx}
\usepackage{booktabs}
\usepackage{array}
\usepackage{multirow}
\usepackage{multicol}

% Listen und Nummerierung
\usepackage{enumitem}
\setlist[itemize]{leftmargin=2em}
\setlist[enumerate]{leftmargin=2em}

% Zeilenabstand
\usepackage{setspace}
\onehalfspacing

% Kopf- und Fußzeilen
\usepackage{scrlayer-scrpage}
\pagestyle{scrheadings}
\clearscrheadfoot

% Code und Quelltext
\usepackage{listings}
\usepackage{verbatim}

\lstset{
    basicstyle=\small\ttfamily,
    breaklines=true,
    breakatwhitespace=true,
    language=Java,
    numbers=left,
    numberstyle=\tiny,
    stepnumber=1,
    showstringspaces=false,
    tabsize=4,
    backgroundcolor=\color{gray!5},
    frame=tb,
    captionpos=b
}

% Mathematik
\usepackage{amsmath}
\usepackage{amssymb}
\usepackage{mathtools}

% Links und Hyperlinks
\usepackage[colorlinks=true, linkcolor=blue, citecolor=blue, urlcolor=blue]{hyperref}

% Abkürzungen
\usepackage[acronym,toc]{glossaries}
\makeglossaries

% Weitere nützliche Pakete
\usepackage{booktabs}
\usepackage{float}
\usepackage[justification=centering]{caption}
\usepackage{subcaption}

%% ============================================================================
%% BEFEHLE
%% ============================================================================

% Blendet Todos aus/ein
\newcommand{\todo}[1]{\textcolor{red}{[TODO: #1]}}

% Abkürzung mit Kürzel
\newcommand{\ac}[1]{\gls{#1}}

%% PDFMetadaten
\hypersetup{
    pdftitle={\doctitle},
    pdfauthor={\authorname},
    pdfsubject={Bachelorarbeit},
}
